% =========================================================
% Rational Numbers
% =========================================================

\subsection{Rational Numbers}

\subsubsection{Definitions and Theorems}

\begin{definition}[Fraction Pairs]
\label{def:fraction-pairs}

Let $\mathbb{Z}$ denote the set of integers.

Define
\[
W = \mathbb{Z} \times (\mathbb{Z} \setminus \{0\})
\]

Thus

\[
W = \{(a,b) : a,b \in \mathbb{Z},\ b \ne 0\}.
\]

Each ordered pair $(a,b)$ is intended to represent a fraction
\[
\frac{a}{b}.
\]

\end{definition}

\begin{remark}[Numerator and Denominator]
\label{rem:numerator-denominator}

For an ordered pair $(a,b) \in W$ with $b \neq 0$, the integer $a$ is called the
\emph{numerator} and the integer $b$ is called the \emph{denominator}.

The pair $(a,b)$ is intended to represent the rational number
\[
\frac{a}{b}.
\]

\end{remark}
L
\begin{remark}
Different pairs may represent the same rational number.
For example

\[
(1,2), (2,4), (3,6)
\]

should represent the same number.  
To formalize this, we introduce an equivalence relation on $W$.

% ---------------------------------------------------------
% Equality of Fraction Pairs
% ---------------------------------------------------------

\begin{definition}[Equality of Fraction Pairs]
\label{def:fraction-pair-equality}

Let
\[
W = \{(a,b) : a,b \in \mathbb{Z},\ b \neq 0\}.
\]

Define a relation $=$ on $W$ by

\[
(a,b) = (c,d)
\quad \Longleftrightarrow \quad
ad = bc .
\]

In words, two ordered pairs represent the same rational number
precisely when the cross products of their components are equal.

\medskip

\noindent
\textbf{Fully quantified form}

\[
\forall (a,b),(c,d) \in W,
\quad
\bigl((a,b) = (c,d)\bigr)
\;\Longleftrightarrow\;
(ad = bc).
\]

\end{definition}

\begin{remark}
This relation identifies different fraction pairs that represent the
same rational number. For example,

\[
(1,2) = (2,4)
\]

because

\[
1 \cdot 4 = 2 \cdot 2.
\]

The rational numbers will be defined as the equivalence classes of
this relation.
\end{remark}

% ---------------------------------------------------------
% Equivalence Relation on Fraction Pairs
% ---------------------------------------------------------

\begin{theorem}
\label{thm:fraction-pair-equivalence}

Let
\[
W = \{(a,b) : a,b \in \mathbb{Z},\ b \neq 0\}.
\]

Define a relation on $W$ by

\[
(a,b) \sim (c,d)
\quad \Longleftrightarrow \quad
ad = bc .
\]

Then $\sim$ is an equivalence relation on $W$. Consequently, the
equivalence classes of $\sim$ partition $W$ and correspond to the
rational numbers.

\end{theorem}

\begin{remark}
The rational numbers are defined as the set of equivalence classes
\[
\mathbb{Q} = W / \sim .
\]
Each equivalence class $[(a,b)]$ represents the rational number
$\frac{a}{b}$.
\end{remark}


% ---------------------------------------------------------
% Rational Numbers
% ---------------------------------------------------------

\begin{definition}[Rational Numbers]
\label{def:rational-numbers}

Let
\[
W = \{(a,b) : a,b \in \mathbb{Z},\ b \neq 0\}
\]

and define the relation

\[
(a,b) \sim (c,d) \iff ad = bc .
\]

The set of \emph{rational numbers} is defined as the set of
equivalence classes of $W$ under this relation:

\[
\mathbb{Q} = W / \sim .
\]

For $(a,b) \in W$, the equivalence class is denoted

\[
[(a,b)] .
\]

We interpret this class as the rational number

\[
\frac{a}{b}.
\]

\end{definition}


% ---------------------------------------------------------
% Distinguished Rational Numbers
% ---------------------------------------------------------

\begin{definition}[Distinguished Elements of $\mathbb{Q}$]
\label{def:rational-distinguished}

The rational number zero is defined as

\[
0 := [(0,1)] .
\]

The rational number one is defined as

\[
1 := [(1,1)] .
\]

\end{definition}


\begin{remark}
The definitions of $0$ and $1$ are well-defined because

\[
(0,1) \sim (0,b)
\quad\text{for all } b \neq 0,
\]

and

\[
(1,1) \sim (a,a)
\quad\text{for all } a \neq 0.
\]

Thus each represents a unique equivalence class.
\end{remark}

% ---------------------------------------------------------
% Addition and Multiplication in Q
% ---------------------------------------------------------

\begin{definition}[Addition and Multiplication in $\mathbb{Q}$]
\label{def:rational-operations}

Let
\[
[(a,b)], [(c,d)] \in \mathbb{Q},
\]
where $(a,b),(c,d) \in W$.

Addition is defined by

\[
[(a,b)] + [(c,d)] := [(ad + bc,\, bd)] .
\]

Multiplication is defined by

\[
[(a,b)] \cdot [(c,d)] := [(ac,\, bd)] .
\]

% ---------------------------------------------------------
% Interpreting Rational Numbers
% ---------------------------------------------------------

\begin{remark}[Three Levels of Representation]
\label{rem:rational-three-levels}

When constructing the rational numbers, it is important to distinguish
between three different mathematical objects that are often written
with similar symbols.

\begin{center}
\begin{tabular}{lll}
\toprule
\textbf{Object} & \textbf{Lives in} & \textbf{Meaning} \\
\midrule
$(a,b)$ & $W$ & Raw integer pair \\
$[(a,b)]$ & $\mathbb{Q}$ & Rational number (equivalence class) \\
$\dfrac{a}{b}$ & notation & Informal representation of the class \\
\bottomrule
\end{tabular}
\end{center}

The ordered pair $(a,b)$ is an element of the set

\[
W = \{(a,b) \in \mathbb{Z} \times \mathbb{Z} : b \neq 0\}.
\]

It is merely a pair of integers and is \emph{not} itself a rational
number.

The rational number corresponding to this pair is the equivalence class

\[
[(a,b)] = \{(c,d) \in W : ad = bc\}.
\]

Thus a rational number is a \emph{set of pairs} that all represent the
same ratio.

The expression

\[
\frac{a}{b}
\]

is simply a convenient notation used to denote the equivalence class
$[(a,b)]$.

\end{remark}

\begin{remark}[Why the Distinction Matters]
Although we commonly write rational numbers as fractions, the
construction shows that the true mathematical object is an equivalence
class of integer pairs.

For example,

\[
\frac{1}{2} = \frac{2}{4} = \frac{3}{6}
\]

because

\[
(1,2) \sim (2,4) \sim (3,6).
\]

All of these pairs belong to the same equivalence class

\[
[(1,2)].
\]

Thus the rational number $\tfrac12$ is not a single pair but the set

\[
[(1,2)] =
\{(1,2),(2,4),(3,6),(4,8),\ldots,(-1,-2),(-2,-4),\ldots\}.
\]

\end{remark}

\begin{remark}[Conceptual Pattern]
This construction is an example of a very common pattern in
mathematics:

\[
\text{set} + \text{equivalence relation} \quad \longrightarrow \quad
\text{quotient structure}.
\]

Here the quotient set

\[
\mathbb{Q} = W / \sim
\]

is obtained by identifying pairs of integers that represent the same
ratio.

Similar constructions appear throughout mathematics, including
modular arithmetic, quotient groups, quotient rings, and quotient
topological spaces.

\end{remark}

\end{definition}


% ---------------------------------------------------------
% Rational Numbers as Infinite Sets
% ---------------------------------------------------------

\begin{remark}[A Rational Number is an Infinite Set]
\label{rem:rational-infinite-set}

By construction, a rational number is an equivalence class of pairs
of integers. That is,

\[
[(a,b)] = \{(c,d) \in W : ad = bc\}.
\]

This equivalence class contains infinitely many elements.

Indeed, if $(a,b) \in W$, then for every integer $k \neq 0$,

\[
(ka,kb) \sim (a,b)
\]

because

\[
(ka)b = k(ab) = a(kb).
\]

Thus

\[
(ka,kb) \in [(a,b)].
\]

Since there are infinitely many integers $k$, the set $[(a,b)]$
contains infinitely many pairs.

Therefore every rational number is an infinite set.

\end{remark}

\begin{remark}[Example]
For example, the rational number $\tfrac12$ is the equivalence class

\[
[(1,2)] =
\{(1,2),(2,4),(3,6),(4,8),\ldots,(-1,-2),(-2,-4),\ldots\}.
\]

All of these pairs represent the same ratio, and therefore correspond
to the same rational number.

\end{remark}


% ---------------------------------------------------------
% Why Well-Definedness Matters
% ---------------------------------------------------------

\begin{remark}[From Representatives to Quotient Structures]
\label{rem:representatives-quotients}

The definitions of addition and multiplication on $\mathbb{Q}$ were
given using \emph{representatives} of equivalence classes:

\[
[(a,b)] + [(c,d)] := [(ad+bc,\,bd)], 
\qquad
[(a,b)] \cdot [(c,d)] := [(ac,\,bd)].
\]

At first glance this may seem problematic.  
A rational number is not a pair $(a,b)$ but the entire equivalence
class $[(a,b)]$. Since the same rational number can be represented by
many different pairs, we must ensure that the result of an operation
does not depend on which representative we choose.

For example,

\[
\frac12 = \frac24.
\]

If we compute

\[
\frac12 + \frac13
\]

using the representatives $(1,2)$ and $(1,3)$ we obtain

\[
\frac{5}{6}.
\]

But if we instead use the representatives $(2,4)$ and $(1,3)$ we obtain

\[
\frac{10}{12}.
\]

These results look different, but they represent the same rational
number because

\[
(5,6) \sim (10,12).
\]

The well-definedness proofs show that this always happens: no matter
which representatives are chosen, the resulting equivalence class is
the same.

\medskip

In general, this situation occurs whenever we construct a new
mathematical object as a set of equivalence classes.  Operations are
often defined on the \emph{representatives}, and we must prove that the
result does not depend on the choice of representative.

Symbolically, the pattern is

\[
\text{representatives}
\quad \longrightarrow \quad
\text{operations}
\quad \longrightarrow \quad
\text{well-defined operations on equivalence classes}.
\]

Once this is established, the quotient set becomes a genuine algebraic
structure.

\end{remark}

\begin{remark}[A Common Pattern in Mathematics]
\label{rem:quotient-pattern}

The construction of the rational numbers is the first example of a
general principle that appears throughout mathematics:

\[
\text{set} + \text{equivalence relation}
\quad \longrightarrow \quad
\text{quotient structure}.
\]

Operations are first defined on representatives and then shown to be
well-defined on equivalence classes.

This idea appears in many important constructions:

\begin{center}
\begin{tabular}{lll}
\toprule
\textbf{Structure} & \textbf{Representatives} & \textbf{Resulting Quotient} \\
\midrule
$\mathbb{Q}$ & integer pairs $(a,b)$ & rational numbers \\
$\mathbb{Z}/n\mathbb{Z}$ & integers & integers modulo $n$ \\
quotient groups & group elements & cosets $G/N$ \\
quotient rings & ring elements & factor rings \\
projective geometry & vectors & projective points \\
\bottomrule
\end{tabular}
\end{center}

Thus the well-definedness proofs for the rational numbers provide the
first example of what is sometimes called \emph{quotient algebra}: the
construction of new algebraic structures from equivalence classes.

\end{remark}



\begin{remark}[Canonical Representative]
Every rational number has a unique representation

\[
\frac{a}{b}
\]

with $\gcd(a,b)=1$ and $b>0$.

This pair is called the \emph{reduced form} of the rational number.
\end{remark}

% =========================================================
% Structural Consequences of the Quotient Construction
% =========================================================

\subsection{Structural Consequences of the Construction of $\mathbb{Q}$}

\begin{remark}[The Partition Structure of $\mathbb{Q}$]
\label{rem:q-partition-structure}

Recall that

\[
W = \{(a,b)\in\mathbb{Z}\times\mathbb{Z} : b\neq0\}
\]

and that the relation

\[
(a,b) \sim (c,d) \iff ad = bc
\]

is an equivalence relation on $W$.

The set of rational numbers is defined as the quotient

\[
\mathbb{Q} = W/\sim .
\]

Thus the rational numbers are the equivalence classes that partition
the set $W$.

Each rational number is therefore a subset of $W$ consisting of all
integer pairs that represent the same ratio.

\end{remark}

% ---------------------------------------------------------

\begin{remark}[A Rational Number is an Infinite Set]
\label{rem:rational-infinite-set}

Let $(a,b)\in W$. The rational number represented by this pair is the
equivalence class

\[
[(a,b)] = \{(c,d)\in W : ad = bc\}.
\]

For every integer $k\neq0$,

\[
(ka,kb) \sim (a,b).
\]

Thus

\[
(ka,kb) \in [(a,b)].
\]

Since there are infinitely many integers $k$, the equivalence class
$[(a,b)]$ contains infinitely many elements.

Therefore every rational number is an infinite set.

\end{remark}

% ---------------------------------------------------------

\begin{remark}[Canonical Representatives]
\label{rem:reduced-rational}

Although a rational number has infinitely many representatives,
there is a natural choice of a preferred representative.

Every rational number can be written uniquely in the form

\[
\frac{a}{b}
\]

where

\[
\gcd(a,b)=1
\qquad
b>0.
\]

This representation is called the \emph{reduced form} of the rational
number.

\end{remark}

% ---------------------------------------------------------

\begin{remark}[Integers Inside $\mathbb{Q}$]
\label{rem:z-in-q}

The integers embed naturally into the rational numbers.

Define the map

\[
n \mapsto [(n,1)].
\]

Thus

\[
\mathbb{Z} \subseteq \mathbb{Q}.
\]

The equivalence class

\[
[(n,1)]
\]

contains all pairs

\[
(n,1),(2n,2),(3n,3),\ldots
\]

which all represent the same integer.

\end{remark}

% ---------------------------------------------------------

\begin{remark}[Geometric Interpretation]
\label{rem:rational-geometry}

Pairs $(a,b)$ can be viewed as points in the integer lattice
$\mathbb{Z}^2$.

Two pairs represent the same rational number precisely when they lie
on the same line through the origin.

Indeed,

\[
(a,b) \sim (c,d)
\]

means

\[
ad = bc,
\]

which is equivalent to

\[
\frac{a}{b} = \frac{c}{d}.
\]

Thus each rational number corresponds to a line through the origin
with rational slope.

\end{remark}

% ---------------------------------------------------------

\begin{remark}[Countability of $\mathbb{Q}$]
\label{rem:q-countable}

Since

\[
W \subseteq \mathbb{Z}^2
\]

and $\mathbb{Z}^2$ is countable, the quotient set

\[
\mathbb{Q} = W/\sim
\]

is also countable.

Thus there exists a bijection

\[
\mathbb{N} \longleftrightarrow \mathbb{Q}.
\]

\end{remark}

% ---------------------------------------------------------

\begin{remark}[A Common Mathematical Pattern]
\label{rem:quotient-pattern}

The construction of the rational numbers is the first example of a
general mathematical pattern:

\[
\text{set} + \text{equivalence relation}
\quad\longrightarrow\quad
\text{quotient structure}.
\]

Operations are defined on representatives and then shown to be
well-defined on equivalence classes.

This idea appears throughout mathematics.

\begin{center}
\begin{tabular}{lll}
\toprule
Structure & Representatives & Resulting Quotient \\
\midrule
$\mathbb{Q}$ & integer pairs $(a,b)$ & rational numbers \\
$\mathbb{Z}/n\mathbb{Z}$ & integers & integers modulo $n$ \\
quotient groups & group elements & cosets $G/N$ \\
quotient rings & ring elements & factor rings $R/I$ \\
projective geometry & vectors & projective points \\
\bottomrule
\end{tabular}
\end{center}

The well-definedness proofs for the operations on $\mathbb{Q}$ are
therefore the first example of what is often called
\emph{quotient algebra}.

\end{remark}

% ---------------------------------------------------------

\begin{remark}[Density of the Rational Numbers]
\label{rem:q-mediant}

Between any two distinct rational numbers

\[
\frac{a}{b} < \frac{c}{d}
\]

there exists another rational number

\[
\frac{a+c}{b+d}.
\]

This fraction is called the \emph{mediant}.  It satisfies

\[
\frac{a}{b}
<
\frac{a+c}{b+d}
<
\frac{c}{d}.
\]

Thus the rational numbers are dense in the real line.

\end{remark}

\begin{remark}[Canonical Representatives]
\label{rem:reduced-rational}

Although a rational number has infinitely many representatives,
there is a natural way to choose a preferred representative.

Let $(a,b)\in W$.  Let

\[
g = \gcd(a,b).
\]

Then

\[
(a,b) \sim \left(\frac{a}{g},\frac{b}{g}\right).
\]

The pair

\[
\left(\frac{a}{g},\frac{b}{g}\right)
\]

satisfies

\[
\gcd\!\left(\frac{a}{g},\frac{b}{g}\right)=1.
\]

Thus every rational number has a representative whose numerator and
denominator are relatively prime.

If we also require

\[
b>0,
\]

this representative is unique.

Therefore every rational number can be written uniquely in the form

\[
\frac{a}{b}
\]

with

\[
\gcd(a,b)=1,
\qquad
b>0.
\]

This representation is called the \emph{reduced form} of the rational
number.

\end{remark}





\end{remark}